\chapter{Conclusions}
\label{chap:conclusions}

\section{Summary of Contributions}
\label{sec:contributions-summary}

This project addressed a well-defined problem: engineering students operate across a fragmented ecosystem of disconnected AI tools, and the cognitive cost of this fragmentation directly reduces learning and productivity. Engunity~X is the answer---a unified, context-aware SaaS platform integrating document retrieval, autonomous research, secure code execution, and adversarial reasoning within a single shared-memory environment.

\subsection{OmniRAG Pipeline}
The OmniRAG Pipeline combines a fine-tuned DistilBERT complexity classifier with four distinct retrieval strategies, achieving 88.0\% retrieval accuracy on multi-hop queries---a 24-percentage-point improvement over a standard single-strategy RAG baseline. HyDE query expansion, CRAG confidence-gated web fallback, and extractive contextual compression each contribute measurably to this result.

\subsection{DeepResearchAgent}
The DeepResearchAgent's five-phase architecture (Decompose, Search, Evaluate, Refine, Synthesize) with configurable depth tiers handles research tasks ranging from a two-minute factual lookup to an eight-iteration comprehensive domain survey. The gap-analysis-driven refine phase---in which the agent re-queries under-evidenced sub-questions---is, to our knowledge, a novel contribution to the open-source agent literature.

\subsection{Multimodal Input Layer}
The integration of CLIP, YOLOv8n, and EasyOCR into a unified \texttt{MultiModalVisionProcessor} enables Engunity to handle engineering image inputs without a separate processing pipeline. Projecting CLIP embeddings into the same FAISS index used for text enables cross-modal retrieval without additional infrastructure.

\subsection{Secure Code Laboratory}
The Code Lab provides a containerised execution environment with an interactive WebSocket terminal, Monaco Editor IDE, DAP-compliant Python debugger, and bidirectional AI error-correction loop. The error correction loop achieves 70.7\% single-iteration success across six error categories, with 91.4\% success on syntax errors.

\subsection{Evidence Quality Score and Decision Vault}
The Decision Vault's EQS metric correlates with human expert ratings at Pearson's $r = 0.81$, correctly flags 78\% of reasoning errors, and maintains a false-positive rate of 11\%. This directly addresses the sycophancy problem in RLHF-trained models.

\section{Objective Review}
\label{sec:objective-review}

\begin{table}[ht]
    \centering
    \caption{Review of project objectives against measured outcomes.}
    \label{tab:objective-review}
    \renewcommand{\arraystretch}{1.3}
    \begin{tabular}{p{4cm} p{3.5cm} p{2.5cm}}
        \toprule
        \textbf{Objective} & \textbf{Target} & \textbf{Outcome} \\
        \midrule
        OmniRAG multi-hop accuracy & $\geq 85$\% & \textbf{88.0\%} \checkmark \\
        P95 latency at 500 users & $< 500$\,ms & \textbf{483\,ms} \checkmark \\
        Code error correction rate & Not specified & \textbf{70.7\%} \\
        EQS correlation with human & Not specified & \textbf{$r=0.81$} \\
        EQS reasoning error detection & Not specified & \textbf{78\%} \checkmark \\
        Contextual compression & $\geq 60$\% token reduction & \textbf{67.7\%} \checkmark \\
        Pilot study satisfaction & Not specified & \textbf{4.25/5.0} \\
        \bottomrule
    \end{tabular}
\end{table}

All specified quantitative targets were met or exceeded. The latency target (P95 $<$ 500\,ms at 500 concurrent users) was met at 483\,ms. The retrieval accuracy target of 85\% on multi-hop queries was met at 88\%.

\section{Limitations}
\label{sec:limitations}

\subsection{Scalability Ceiling}
The single-node deployment with in-process FAISS indexing limits horizontal scalability. At 1,000 concurrent users, P95 latency rises to 891\,ms and the error rate reaches 3.2\%. Addressing this requires migrating to an external vector store (Qdrant or Weaviate) and deploying multiple stateless FastAPI instances behind a load balancer.

\subsection{Knowledge Graph Staleness}
Community summaries are regenerated on a 6-hour schedule. Documents ingested between worker runs are retrievable via vector search but are not reflected in community summaries used by the GraphRAG path. Incremental graph updates would reduce this staleness without the cost of full recomputation.

\subsection{YOLOv8n Taxonomy Mismatch}
YOLOv8n was trained on COCO 2017, which includes 80 everyday object classes but no engineering-specific hardware. Engineering-specific objects are detected as the closest COCO proxy with degraded confidence. A fine-tuned YOLOv8 model on an engineering hardware dataset would substantially improve utility.

\subsection{Session Length Limit}
The hierarchical session summary compression is effective for sessions up to approximately 100 turns. Beyond this, the compressed summary approaches the context window limit. A sliding window summary approach would extend the practical session length limit.

\section{Future Work}
\label{sec:future-work}

Six concrete directions are identified for future work:
\begin{enumerate}
    \item \textbf{External Vector Store}: Migrate from in-process FAISS to Qdrant to enable horizontal scaling of the retrieval layer beyond the current single-node ceiling.
    \item \textbf{Engineering-Domain Fine-Tuning}: Fine-tune YOLOv8 on an annotated engineering hardware dataset; fine-tune DistilBERT on engineering-domain queries; fine-tune a text embedding model on engineering textbooks.
    \item \textbf{Incremental Knowledge Graph Updates}: Replace full-graph recomputation with incremental entity and edge insertion, reducing community staleness from 6 hours to near-real-time.
    \item \textbf{Formal User Study}: A randomised controlled trial comparing Engunity and control cohorts on a shared final-year project task, measuring time-to-completion, solution quality, and NASA-TLX cognitive load.
    \item \textbf{Test-Case-Augmented Error Correction}: Provide failing test cases alongside error messages to improve logic-error correction rate from the current 47\% toward the 85\%+ target.
    \item \textbf{Mobile-Responsive PWA}: A mobile-responsive redesign with PWA manifest enabling offline caching, extending accessibility to students working on mobile devices.
\end{enumerate}

\section{Ethical Considerations}
\label{sec:ethics}

All user data is scoped via Row-Level Security policies in PostgreSQL and Supabase. Users can delete their data at any time via the Settings interface, triggering a cascading delete across all related tables. No user data is shared with third-party services except through Groq API and Brave Search API, invoked with keys that do not include user-identifying information.

Engunity's AI-generated responses are accompanied by self-critique confidence scores. The Decision Vault's adversarial challenges are framed as questions rather than refutations, encouraging critical thinking. Like any AI tool, Engunity can be misused for academic dishonesty; its deployment in educational settings should be accompanied by clear institutional guidelines on appropriate use.

\section{Final Remarks}
\label{sec:final-remarks}

Engunity~X demonstrates that high-performance, production-quality AI infrastructure for engineering education is achievable using entirely open-source components on commodity hardware. The five-service architecture deploys in under 30 minutes on any Linux host with Docker. The measured results---88\% multi-hop retrieval accuracy, 483\,ms P95 latency at 500 users, 78\% reasoning-error detection, 4.25/5.0 pilot satisfaction---confirm the design achieves its intended goals. The architecture's most important quality is its \textit{traceability}: every pipeline stage has defined inputs, outputs, and failure modes, every design decision has a documented rationale, and every performance claim is backed by reproducible experiments.
